% ---- RANDOM FORMATTING ---- %
\setlength{\parindent}{0pt}
%\renewcommand{\thesection}{\Roman{section}} 


% ---- PACKAGES ---- %
\usepackage[margin=2cm]{geometry}
\usepackage{tikz} % graphs
\usepackage{tikz-cd} % commutative diagrams
\usetikzlibrary{decorations.pathreplacing} % for curly braces 
\usetikzlibrary{arrows,shapes,positioning,calc}
\usetikzlibrary{decorations.markings}
\usepackage{amsmath} % ams lot
\usepackage{amsthm} 
\usepackage{amssymb}
\usepackage{parallel}
\usepackage[normalem]{ulem} % this is for a stupid joke in my Galois representations notes
\usepackage{bbm} % blackboard bold numbers
%\usepackage{mathpazo} % for blackboard bold
%\usepackage{amsfonts}
%\usepackage{mathabx}
%\usepackage{MnSymbol}
%\usepackage{newpxtext} % font
%\usepackage{tgpagella} % font
\usepackage{epigrafica} % font
\usepackage{fourier} % font
%\usepackage{mathptmx} % maths font
%\usepackage{txfonts} % font
\usepackage[bb=ams, cal=cm, scr=rsfs, bbsymbols, calscaled=.96, bbscaled=.96]{mathalpha} % sets the calligraphic to mathcal, scr to rsfs, and blackboard bold to ams
\NewCommandCopy{\pazomathbb}{\mathbb} % if using fourier, this and the line above stop it from using the horrible looking fourier bb font
\AtBeginDocument{\RenewCommandCopy{\mathbb}{\pazomathbb}} % ^
\usepackage{enumerate} % fancy labelling for lists 
% \usepackage{enumitem} % resuming lists - CLASHES WITH ABOVE
\usepackage{pgfplots}
\pgfplotsset{compat=1.18}
\usepgfplotslibrary{fillbetween}
\usepackage{float}
\usepackage{caption} % i do not entirely know -- some latex on stackexchange taught me how to do side-by-side figures
\usepackage{subcaption}
\usepackage[dvipsnames]{xcolor}
\usepackage[most]{tcolorbox} % for fancy boxes
\tcbuselibrary{skins,breakable,theorems} % this is for fancy box formatting
\usepackage{fontawesome}
%\usepackage[sc]{titlesec} % smallcaps section headings
\usepackage{imakeidx} % index generator
\makeindex

\usepackage[hidelinks]{hyperref} % clicky links

% ---- TIKZ ---- %
\usepackage{tikz}

\usetikzlibrary{arrows,shapes,positioning}
\usetikzlibrary{decorations.markings}
\usetikzlibrary{fadings,patterns}

% SNAKE ARROWS STYLE! The argument is the angle
\tikzcdset{snaky/.style={rounded corners,to path={[rotate=#1] -- ($(\tikztostart.#1)+(2ex,0)$)%
 -- ({$(\tikztostart.#1)+(2ex,0)$} |- {$0.5*(\tikztostart)+0.5*(\tikztotarget)$})%
 -- ({$(\tikztotarget.\number\numexpr 180+(#1) \relax)+(-2ex,0)$} |- {$0.5*(\tikztostart)+0.5*(\tikztotarget)$}) \tikztonodes%
 -- ($(\tikztotarget.{\number\numexpr 180+(#1) \relax})+(-2ex,0)$)%
 -- (\tikztotarget)}}}


\tikzstyle arrowstyle=[scale=1]
\tikzstyle directed=[postaction={decorate,decoration={markings, mark=at position .6 with {\arrow[arrowstyle]{stealth}}}}]
\tikzstyle Node=[circle, draw=black, fill=black!10]

\pgfplotsset{compat=1.18}


% ---- COMMANDS ---- %
% mathbb commands
\newcommand{\A}{\mathbb{A}}
\newcommand{\D}{\mathbb{D}}
\newcommand{\R}{\mathbb{R}}
\newcommand{\Q}{\mathbb{Q}}
\newcommand{\Pruf}{\mathbb{D}}
\newcommand{\Z}{\mathbb{Z}}
\newcommand{\N}{\mathbb{N}}
\newcommand{\F}{\mathbb{F}}
\newcommand{\bbG}{\mathbb{G}}
\newcommand{\PP}{\mathbb{P}}
\newcommand{\Hyp}{\mathbb{H}}
\newcommand{\Uni}{\mathbb{U}\mathrm{ni}}
\newcommand{\Gm}{\bbG_{\mathrm{m}}}
\newcommand{\Ga}{\bbG_{\mathrm{a}}}
\newcommand{\triv}{\mathbbm{1}}

% mathcal/mathscr commands
\newcommand{\E}{\mathcal{E}}
\newcommand{\Hopf}{\mathscr{H}}
\newcommand{\GHopf}{\mathscr{G}}
\newcommand{\GLHopf}{\mathscr{G\!l}\!}
\newcommand{\SLHopf}{\mathscr{S\!l}\!}
\newcommand{\inducedop}[2][A]{{#2}_{|_{#1}}}
\newcommand{\ints}{\mathcal{O}}
\newcommand{\Autforms}{\mathcal{A}}
\newcommand{\modforms}{\mathcal{M}}
\newcommand{\cuspforms}{\mathcal{S}}
\newcommand{\Ell}{\mathcal{L}}
\newcommand{\Hilb}{\mathcal{H}}
\newcommand{\calC}{\mathcal{C}}
\newcommand{\calD}{\mathcal{D}}
\newcommand{\calQ}{\mathcal{Q}}
\newcommand{\scrR}{\mathscr{R}}

% mathbf/boldsymbol commands
\newcommand{\Set}{\mathbf{Set}}         % (Set)s
\newcommand{\Grp}{\mathbf{Grp}}         % (Gr)ou(p)s
\newcommand{\Vect}{\mathbf{Vect}}       % (Vect)or spaces
\newcommand{\Mod}{\mathbf{Mod}}         % (Mod)ules
\newcommand{\CRing}{\mathbf{CRing}}     % (C)ommutative (Ring)s
\newcommand{\CAlg}{\mathbf{CAlg}}		% (C)ommutative (Alg)ebras
\newcommand{\Sep}{\mathbf{Sep}}         % (Sep)arable field extensions
\newcommand{\bfrho}{\boldsymbol{\rho}}  % universal deformation
\newcommand{\bfchi}{\boldsymbol{\chi}}  % universal character


% mathrm commands
\newcommand{\Gal}{\mathrm{Gal}} % (Gal)ois group
\newcommand{\Sel}{\mathrm{Sel}} % (Sel)mer group
\newcommand{\Hom}{\mathrm{Hom}} % (Hom)set
\newcommand{\End}{\mathrm{End}} % (End)omorphisms
\newcommand{\Aut}{\mathrm{Aut}} % (Aut)omorphisms
\newcommand{\Frob}{\mathrm{Frob}} % (Frob)enius elements
\newcommand{\Frat}{\mathrm{Frat}} % (Frat)tini quotient
\newcommand{\Rep}{\mathrm{Rep}} % category of (Rep)resentations
\newcommand{\Def}{\mathrm{Def}} % (Def)ormation functor
\newcommand{\DefQ}{\Def_{\calQ}} % (Def)ormation functor with property (Q)
\newcommand{\Irr}{\mathrm{Irr}} % (Irr)educibles
\newcommand{\Isom}{\mathrm{Isom}} % (Isom)etries
\newcommand{\Ext}{\mathrm{Ext}} % (Ext)ensions
\newcommand{\Cl}{\mathrm{Cl}} % (Cl)ass group
\newcommand{\K}{\mathrm{K}} % K theory
\newcommand{\cohom}{\mathrm{H}} % co(H)omology
\newcommand{\res}{\mathrm{res}} % (res)triction
\newcommand{\Res}{\mathrm{Res}} % (Res)triction -- capitalised for rep theory
\newcommand{\cores}{\mathrm{cores}} % (cores)triction
\newcommand{\Ind}{\mathrm{Ind}} % (Ind)uction
\newcommand{\ab}{\mathrm{ab}} % (ab)elianisation
\newcommand{\rad}{\mathrm{R}} % (R)adical (algebraic groups)
\newcommand{\unip}[1]{{#1}_{\mathrm{u}}} % (u)nipotent (subscript)
\newcommand{\uniprad}{\mathrm{R}_{\mathrm{u}}} % unipotent radical
\newcommand{\Pgp}{\mathrm{P}} % P (projection to (P)GL)
\newcommand{\Special}{\mathrm{S}} % S (intersection with (S)L)
\newcommand{\GL}{\mathrm{GL}} % classical groups below
\newcommand{\SL}{\mathrm{SL}}
\newcommand{\PSL}{\Pgp\SL}
\renewcommand{\O}{\mathrm{O}}
\newcommand{\SO}{\mathrm{SO}}
\newcommand{\SU}{\mathrm{SU}}
\newcommand{\Sp}{\mathrm{Sp}}
\newcommand{\Gr}{\mathrm{Gr}} % (Gr)assmannian
\newcommand{\Spec}{\mathrm{Spec}} % Affine scheme
\newcommand{\Spf}{\mathrm{Spf}} % Formal scheme
\newcommand{\Proj}{\mathrm{Proj}} % (Proj)ective affine scheme
\newcommand{\Frac}{\mathrm{Frac}} % Field of (Frac)tions
\newcommand{\ev}{\mathrm{ev}} % (ev)aluation map
\newcommand{\coker}{\mathrm{coker}} % (coker)nel
\newcommand{\im}{\mathrm{im}} % (im)age
\newcommand{\tr}{\mathrm{tr}} % (tr)ace
\newcommand{\lcm}{\mathrm{lcm}} % (l)east (c)ommon (m)ultiple
\newcommand{\trans}{\mathrm{trf}} % (tr)ans(f)er map
\newcommand{\norm}{\mathrm{N}} % (N)orm map
\newcommand{\milnor}{\mathrm{M}} % (M)ilnor (superscript for K-theory)
\newcommand{\nr}{\mathrm{nr}} % u(nr)amified
\newcommand{\et}{\mathrm{\acute{e}t}} % (ét)ale
\newcommand{\Char}{\mathrm{char}} % (char)acteristic of a ring
\newcommand{\Ad}{\mathrm{Ad}} % (Ad)joint representation


% mathfrak commands
\newcommand{\ratfuncs}{\mathfrak{K}}
\newcommand{\Liealg}{\mathfrak{g}}
\newcommand{\eideal}{\mathfrak{e}}
\newcommand{\hyperb}{\mathfrak{H}}
\newcommand{\prid}{\mathfrak{p}}
\newcommand{\mxid}{\mathfrak{m}}
\newcommand{\nxid}{\mathfrak{n}}
\newcommand{\qprid}{\mathfrak{q}}
\newcommand{\Prid}{\mathfrak{P}}
\newcommand{\ideal}{\trianglelefteq}
\newcommand{\id}{\mathrm{id}}
\newcommand{\cl}{\mathrm{cl}}


\newcommand{\mono}{\hookrightarrow}
\newcommand{\epi}{\twoheadrightarrow}
\newcommand{\iso}{\stackrel{\sim}{\smash{\longrightarrow}\rule{0pt}{0.4ex}}}
\newcommand{\To}{\Rightarrow}
\newcommand{\ot}{\leftarrow}
\newcommand{\oT}{\Leftarrow}
\newcommand{\actson}{\circlearrowright}
\makeatletter
\newcommand{\xRightarrow}[2][]{\ext@arrow 0359\Rightarrowfill@{#1}{#2}}
\makeatother

\newcommand{\gen}[1]{\left\langle#1\right\rangle}
\newcommand{\ggen}[1]{\left\langle\!\left\langle #1\right\rangle\!\right\rangle}
\newcommand{\stein}[1]{\left\{ #1 \right\}}
\newcommand{\sstein}[1]{\stein{\!\!\stein{ #1 }\!\!}}
\newcommand{\ssstein}[1]{\stein{\!\!\sstein{-,-}\!\!}}
\newcommand{\sqbrack}[1]{\left[#1\right]}
\newcommand{\ssqbrack}[1]{\left\llbracket#1\right\rrbracket}
\newcommand{\abs}[1]{\left| #1 \right|}
\newcommand{\aabs}[1]{\left|\!\left| #1 \right|\!\right|}

\newcommand{\wt}[1]{\widetilde{#1}}
\newcommand{\wh}[1]{\widehat{#1}}
\newcommand{\ol}[1]{\overline{#1}}
\newcommand{\ul}[1]{\underline{#1}}
\newcommand{\bemph}[1]{\textcolor{RubineRed}{\emph{#1}}}
\newcommand{\smallfrac}[2]{{\textstyle \frac{#1}{#2}}}
\newcommand{\legendre}[2]{\genfrac{(}{)}{}{}{#1}{#2}}


\newcommand{\eps}{\varepsilon}
\newcommand{\quickmatrix}[1]{\begin{pmatrix} #1 \end{pmatrix}}

\newcommand{\congr}{\mathrm{cg}}
\newcommand{\cgsubgp}{\leq_{\congr}}

\newcommand{\ch}{\mathrm{ch}}

\newcommand{\BS}{\mathrm{BS}} % Borel-Serre compact'n
\newcommand{\Exc}{\textcolor{RubineRed}{$\ast$} \textbf{Exc}} % exercise marker

% Defines \Sha
\usepackage[T2A,T1]{fontenc} % the latter two options are for epigrafica (LGR,OT1) -- IF USING FOURIER PACKAGE, DON'T USE OT1
\DeclareSymbolFont{cyrillic}{T2A}{cmr}{m}{n}
\DeclareMathSymbol{\Sha}{\mathalpha}{cyrillic}{216}

% this bit does xleftrightarrow
\makeatletter
\newcommand\xleftrightarrow[2][]{%
  \ext@arrow 9999{\longleftrightarrowfill@}{#1}{#2}}
\newcommand\longleftrightarrowfill@{%
  \arrowfill@\leftarrow\relbar\rightarrow}
\makeatother


\renewcommand{\C}{\mathbb{C}} % these have to be here because fontenc fucks with them
\renewcommand{\U}{\mathrm{U}}
\renewcommand{\mod}[1]{\,\,\mathrm{mod}\,\, #1}

% Source - https://tex.stackexchange.com/a
% Posted by AboAmmar, modified by community. See post 'Timeline' for change history
% Retrieved 2026-01-14, License - CC BY-SA 4.0

%\numberwithin{equation}{section}


% THIS IS SOME LATEX MAGIC THAT LETS YOU REFER TO CUSTOM ITEMIZE LABELS USING HYPERREF
% to use, use \item[\nameditem{ref:erence}{LABEL}] in an itemize environment
% this renders as `LABEL blah blah blah' and is referenced by \ref{ref:erence}
\makeatletter
\def\nameditem#1#2{\begingroup
	#2%
	\def\@currentlabel{#2}%
	\phantomsection\label{#1}\endgroup
}
% Source - https://tex.stackexchange.com/a/730517/432002, Yerasito


% ---- AMSTHM ENVs ---- %
\newcommand{\newmarkedtheorem}[1]{%
  \newenvironment{#1}
	{%
	  \renewcommand{\qedsymbol}{\ensuremath{\diamondsuit}} % this makes the end symbol of an eg into a diamond
	  \pushQED{\qed}\csname inner@#1\endcsname
	}
	{\popQED\csname endinner@#1\endcsname}%
  \newtheorem{inner@#1}%
}

\theoremstyle{definition}
\newtheorem{DEFN}{Definition}[section]
\newtheorem{prop}[DEFN]{Proposition}
\newtheorem{lemma}[DEFN]{Lemma}
\newmarkedtheorem{eg}[DEFN]{Example}
\newmarkedtheorem{noneg}[DEFN]{Non-example}
\newtheorem*{rmk}{Remark}
\newtheorem{cor}{Corollary}[DEFN]
\newtheorem*{redcor}{\textcolor{RubineRed}{Corollary}}

\theoremstyle{plain}
\newtheorem{THM}[DEFN]{Theorem}
%\newtheorem{conj}[DEFN]{Conjecture}
\newtheorem{ax}{Axiom}

% fancy definitions
\newtcolorbox{DEFNbox}{
	sharp corners,
	breakable,
	top=.35cm,
	bottom=.25cm,
	left=.35cm,
	right=.35cm,
	colback=black!5!white,
	boxrule=-1pt,
	leftrule=3pt,
	colframe=blue!40!red,
	before skip=\topsep,
	after skip=\topsep,
	fonttitle=\large\bfseries\sffamily,
	enhanced,
}

\newtcolorbox{THMbox}{
	sharp corners,
	breakable,
	top=.35cm,
	bottom=.25cm,
	left=.35cm,
	right=.35cm,
	colback=black!5!white,
	boxrule=-1pt,
	leftrule=3pt,
	colframe=RubineRed,
	before skip=\topsep,
	after skip=\topsep,
	fonttitle=\large\bfseries\sffamily,
	enhanced,
}

\newtcolorbox{conj}[1]{
	sharp corners,
	breakable,
	boxrule=-1pt,
	colframe=blue!40!white,
	top=.25cm,
    bottom=.25cm,
    left=.35cm,
    right=.35cm,
	leftrule=3pt,
	title=\faQuestionCircleO \,\,\emph{#1},
    fonttitle=\large\bfseries\sffamily,
	enhanced,
	coltitle=white,
	colback=black!5!white,
	attach boxed title to top left={xshift=1pt, yshift=-.1cm},
	boxed title style={colback=blue!40!white, boxrule=-1pt, sharp corners}
}

\newenvironment{defn}{\leavevmode\begin{DEFNbox}\begin{DEFN}\vspace{-.15cm}}{\end{DEFN}\end{DEFNbox}\leavevmode}
\newenvironment{thm}[1][]{\leavevmode\begin{THMbox}\begin{THM}[#1]\vspace{-.15cm}}{\end{THM}\end{THMbox}\leavevmode}

\DeclareMathVersion{normal2} % this is to fix a "too many math alphabets" problem